\chapter{Conclusões}
\label{chap:conclusoes}

Este trabalho teve como objetivo investigar a aplicação de técnicas de modelagem baseada em dados e de modelagem híbrida, com incorporação de informação física, para a predição de variáveis de desempenho em sistemas de destilação por membranas do tipo V-AGMD. O estudo concentrou-se na previsão do fluxo de permeado e das temperaturas de saída das correntes de alimentação e resfriamento a partir de variáveis operacionais do sistema, por meio de um \textit{pipeline} que integrou análise exploratória dos dados, validação cruzada com separação por regimes experimentais, comparação entre diferentes famílias de modelos e uso de predições oriundas de um modelo físico reduzido do processo.

\subsection{Conclusões gerais e interpretação dos resultados}

A partir das métricas observadas, verificou-se que o melhor desempenho foi obtido pelo modelo híbrido HRNN, isto é, uma rede neural multicamada com incorporação das predições do modelo físico 0D e de restrições baseadas no comportamento fenomenológico do sistema. Esse modelo apresentou redução do erro quadrático médio da raiz (RMSE) em relação ao melhor modelo linear e ao melhor modelo baseado em árvores.

A performance competitiva das famílias de modelos lineares sugere que, na faixa operacional analisada, uma parcela relevante da variabilidade dos alvos pode ser representada por relações de menor complexidade. Ainda assim, o ganho obtido pelo modelo híbrido indica que a representação não linear da rede neural, combinada à informação física fornecida pelo modelo 0D, permitiu capturar efeitos adicionais associados ao acoplamento entre transferência de calor e massa no processo.

Do ponto de vista conceitual, esse resultado pode ser interpretado como uma aproximação de um cenário ideal de modelagem híbrida. O modelo mais adequado para representar o comportamento do sistema seria aquele capaz de reproduzir, de forma eficiente, as equações de balanço de massa e energia e, simultaneamente, incorporar as imperfeições observadas na operação real. Nesse sentido, a estratégia adotada se aproxima desse objetivo ao combinar as predições do modelo físico reduzido com uma rede neural treinada para aprender correções sistemáticas em relação aos dados experimentais.

Sob essa perspectiva, a rede neural fornece uma estrutura capaz de aproximar relações não lineares complexas, enquanto o modelo físico 0D representa, de forma simplificada ou média, os principais acoplamentos fenomenológicos do sistema. A combinação dessas duas abordagens permite que o modelo híbrido incorpore tanto a informação empírica presente nos dados quanto uma tendência física previamente estimada pelo modelo fenomenológico. No entanto, essa aproximação permanece limitada à região de operação contemplada pelos dados experimentais. Assim, mesmo com a incorporação de informação física, a capacidade de generalização do modelo continua dependente da cobertura dos dados disponíveis e não deve ser interpretada como garantia de extrapolação para condições operacionais não avaliadas experimentalmente.

Essa interpretação é particularmente relevante porque o modelo físico reduzido está sujeito a simplificações inerentes à sua formulação. Assim, a rede híbrida atua como um corretor sistemático das discrepâncias entre o modelo físico e o comportamento experimental. Essa complementaridade mostrou-se mais eficaz do que o uso isolado de modelos puramente \textit{data-driven} ou puramente \textit{physics-driven}, ao menos dentro da faixa de operação investigada.

A análise também deve considerar os fundamentos estatísticos da capacidade de generalização. Modelos mais flexíveis podem apresentar excelente ajuste ao conjunto de treinamento sem necessariamente reproduzir novos dados com precisão, especialmente na presença de ruído experimental. Nesse trabalho, o uso de validação cruzada e critérios de seleção com penalização de complexidade foi fundamental para reduzir o risco de sobreajuste. O bom desempenho do modelo HRNN sugere que sua maior flexibilidade foi adequadamente regularizada, permitindo capturar relações consistentes nos dados.


Além disso, as diferentes variáveis de saída estão fisicamente acopladas. O fluxo de permeado depende diretamente dos gradientes térmicos e de pressão de vapor estabelecidos no módulo, bem como das resistências térmicas e difusivas e dos efeitos de polarização. Ao mesmo tempo, as temperaturas de saída também são afetadas pela transferência de massa, uma vez que a evaporação e a condensação do permeado envolvem trocas de calor latente e alteram os balanços de energia nas correntes de alimentação e refrigeração. Assim, embora a seleção de modelos tenha sido orientada prioritariamente pelo erro associado ao fluxo de permeado, a previsão das temperaturas permanece como uma avaliação relevante da capacidade do modelo em representar o comportamento térmico acoplado do sistema.

Os resultados obtidos indicam que, mesmo sem terem sido priorizadas explicitamente no critério de seleção, as temperaturas de saída apresentaram desempenho preditivo elevado. Isso sugere que a arquitetura selecionada foi capaz de capturar, de forma satisfatória, não apenas a variável principal de produtividade, mas também as respostas térmicas associadas ao funcionamento do processo.


Por fim, destaca-se que a estratégia de seleção de modelos foi deliberadamente orientada pela previsão do fluxo de permeado, por este representar o principal indicador de produtividade do processo. Embora todos os modelos tenham sido treinados para prever simultaneamente o fluxo e as temperaturas de saída, a busca e a seleção de hiperparâmetros utilizaram como critério principal o erro associado ao alvo $Flux$. Assim, o mesmo procedimento de priorização não foi repetido separadamente para $Alim\_T\_out$ e $Ref\_T\_out$, a fim de manter uma única estratégia de seleção alinhada ao objetivo central do trabalho e evitar a escolha de modelos distintos para cada variável de saída.

Ainda assim, as temperaturas de saída foram mantidas na avaliação por sua relevância física e energética, uma vez que influenciam diretamente métricas como o consumo específico de energia térmica ($SEC$) e a razão de ganho de saída ($GOR$). Além disso, mesmo não tendo sido priorizadas no critério de seleção, essas variáveis apresentaram desempenho preditivo elevado, indicando que a estratégia adotada foi capaz de preservar uma boa representação das respostas térmicas do sistema. Dessa forma, os resultados obtidos para essas variáveis devem ser interpretados como uma avaliação complementar do modelo selecionado com base no fluxo de permeado.

\subsection{Trabalhos futuros}

A partir dos resultados obtidos, diversas possibilidades de aprofundamento podem ser propostas.

Uma primeira direção consiste em investigar estratégias de modelagem que priorizem explicitamente a predição das temperaturas de saída, por meio de critérios de seleção orientados por erro térmico ou buscas de hiperparâmetros dedicadas a essas variáveis.

Outra linha promissora envolve o refinamento das estratégias de incorporação de conhecimento físico. Embora a abordagem híbrida adotada tenha apresentado bons resultados, ainda é possível explorar outras formas de regularização física, incluindo funções de perda híbridas ou arquiteturas que imponham maior consistência com balanços de energia e massa.

Também pode ser interessante avaliar o uso de \textit{Physics-Informed Neural Networks} (PINNs), especialmente em cenários onde se deseja ampliar a capacidade de generalização para regiões operacionais pouco amostradas. No entanto, para a faixa considerada neste estudo, as estratégias híbridas já se mostraram adequadas.

Uma ampliação natural deste trabalho consiste na obtenção de novos dados experimentais em regiões mais amplas do espaço de operação. Isso permitiria avaliar a robustez dos modelos fora da região atualmente analisada e explorar de forma mais clara possíveis não linearidades do sistema.

Por fim, os modelos desenvolvidos podem ser aplicados em estudos de otimização operacional, permitindo identificar condições de operação que maximizem a produtividade, minimizem o consumo energético ou conciliem ambos os objetivos.

Em síntese, este trabalho mostrou que a integração entre modelos físicos reduzidos e aprendizado de máquina constitui uma abordagem eficaz para a modelagem de sistemas V-AGMD. Os resultados reforçam que a combinação entre conhecimento fenomenológico, validação estatística rigorosa e estratégias de modelagem híbrida é fundamental para obter modelos preditivos robustos, interpretáveis e relevantes do ponto de vista de engenharia.